\section{Coefficient systems and their contraction identities}
\label{co:section}

We construct Fourier coefficients detecting central values and Heegner
points, then use arithmetic lower bounds to normalize their least nonzero
valuation.  The resulting binary tests satisfy the contraction and isolation
identities used in Section~\ref{gr:section}.  The even lower bound is already
available from Section~\ref{cy:section}; the odd detector is constructed
without one, and its normalization remains conditional until
Proposition~\ref{pw:odd-lower}.

Throughout, $S$ contains $2$ and the primes dividing the
conductor $N$.  A filter prescribes local unit squareclasses at $S$, together
with the infinite sign, and excludes finitely many further primes.  Its
indices are $h=\sigma D$, where $D$ is positive and squarefree.  All primes
adjoined subsequently are odd, good, and outside the current support.
Put $F=\mathbb Q(E[2])$.  In the present residual range, $F$ is trivial or
quadratic.  A prime split in $F$ has weight $1$, and a prime inert in $F$,
called simple, has weight $1/2$.  We write $w(D)$ for the sum of these weights.
All valuations satisfy $v(2)=1$.

We use the theta multiplier cubed for weight $3/2$.  Local unit characters
and their associated Dirichlet characters are related by the usual inverse
idele convention.  Thus, if a finite Schwartz function satisfies
$f(ux)=\alpha(u)f(x)$, its scalar classical character on the lower-right
entry is the Dirichlet character associated to $\alpha^{-1}$.
This convention removes a possible inverse ambiguity in the formulas below.

\subsection{Even coefficients and exact comparison of ratios}

\begin{proposition}[Even coefficient system]\label{co:even}
Fix an even-sign filter and a nonvanishing index $D_0$ in it.
There is a holomorphic cuspidal form $g=\sum_{n\geq1}a(n)q^n$ of weight
$3/2$, orthogonal to unary theta series, of level supported on a fixed
enlargement of $S$, and characters $\eta,\psi$ of $2$-power order such that
$\psi/\eta=\chi_{-\sigma}$ and
\begin{equation}\label{co:recurrence}
 a(p^2n)=\bigl(\eta(p)a_p(E)-\psi(p)(-n/p)\bigr)a(n)
              -p\eta(p)^2a(n/p^2)\qquad(p\notin S).
\end{equation}
Its coefficients are zero off the unit-squareclass filter.  After
normalization at $D_0$, every squarefree coefficient is rational up to a
$2$-power root of unity.  For two nonzero squarefree coefficients,
\begin{equation}\label{co:exact-ratio}
 2v\!\left(\frac{a(D)}{a(D_0)}\right)
   =v\!\left(\frac{L_*(\sigma D)}{L_*(\sigma D_0)}\right).
\end{equation}
Moreover $a(D)\neq0$ exactly when $L(E^{(\sigma D)},1)\neq0$.
Consequently any previously established bound
$v(L_*(\sigma D))\geq2w(D)+s(\sigma D)-C$ gives
\begin{equation}\label{co:even-bound}
 v(a(D))\geq w(D)+\tfrac12s(\sigma D)-C_g
\end{equation}
with $C_g$ fixed independently of $D$.
\end{proposition}

\begin{proof}
Use the Shimura packet of $E^{(-\sigma)}$.  Its positive index $D$ tests
the twist by $\chi_{-D}$, which is $E^{(\sigma D)}$.
The nonvanishing at $D_0$ selects a globally occurring member admitting
that Fourier functional.  Take its spherical vector outside $S$ and its
holomorphic vector of weight $3/2$ at infinity.

Here is the support-local vector selection.  Start with a vector on which
the $D_0$ Fourier functional is nonzero.  Projection under a sufficiently
large compact group of upper unipotents to additive frequency $D_0$
restricts the potentially nonzero unit translates to arbitrarily small
neighborhoods of the two solutions of $u^2=1$.  Choose a unit character
agreeing with the forced action at $-1$.  It can have $2$-power order:
the image of $-1$ is detected in the $2$-primary quotient of the unit
group.  The neighborhoods of $1$ and $-1$ then contribute with the same
sign, so this projection preserves the nonzero functional.  Upper
integral translations act trivially, and smoothness gives sufficiently
deep lower-unipotent invariance.  The vector therefore defines a
classical form at a level supported on $S$, allowing arbitrary depth at
$2$ and at the bad primes.  The unit-squareclass and coprimality sieves
commute with good $T(p^2)$, since their conditions are unchanged by
multiplication of the index by $p^2$.  The classical Hecke formula
\cite{Shimura1973HalfIntegral} gives
\eqref{co:recurrence}, with $\eta=\psi\chi_{-\sigma}$.

At the level just chosen, let $V$ be the full space of weight-$3/2$
cusp forms orthogonal to unary theta series, with scalar character
$\psi$ and good $T(p^2)$ eigenvalues $\eta(p)a_p(E)$.  It contains
the constructed form, whose coefficient at $D_0$ is nonzero.  For
$D/D_0$ a unit square at every support prime, Waldspurger's comparison
\cite[Corollary~2, pp.~483--484]{Waldspurger1981} gives, for every
$f\in V$,
\[
 a_f(D)^2=
 \psi(D/D_0)\frac{L_*(\sigma D)}{L_*(\sigma D_0)}a_f(D_0)^2.
\]
Its hypothesis (H1) follows from the nonzero eigenspace
\cite[\S I.2, Proposition~2]{Waldspurger1981}; the final paragraph of
the comparison proof removes (H2) for these local squareclasses.
Factoring the displayed quadratic identity in the two coefficient
functionals shows that they are proportional on all of $V$, with one
scalar independent of $f$.  This includes oldvectors and their linear
combinations; if the central value at $D$ vanishes, the first functional
is simply zero.

Choose an algebraic form $g$ in the filtered subspace of $V$ with
$a(D_0)\ne0$.  To justify this choice, multiplication by ordinary theta
embeds the space into an integral-weight space, and divisibility by
theta is specified by algebraic cusp conditions
\cite[\S\S3--4]{Ramsey2006HalfIntegral}.  The character, good-Hecke
and filter conditions are algebraic linear conditions as well.
Since evaluation at $D_0$ is nonzero after extension to $\mathbb C$,
it is nonzero on some algebraic form.

Let $u$ be an automorphism of $\overline{\mathbb Q}$ acting on the
character values by an odd power $e$, and put $\nu=\eta^{(e-1)/2}$.
The forms $g^u$ and $g\otimes\nu$ have the same scalar character
$\psi^e=\psi\nu^2$ and the same good eigenvalues
$\eta(p)^e a_p(E)=\nu(p)^2\eta(p)a_p(E)$; here $\psi/\eta$ is
quadratic.  They belong to one good-Hecke eigenspace at a common level
still supported on $S$.  Applying the preceding proportionality in
that nonzero space and normalizing at $D_0$ gives
\[
 u(\xi)=\nu(D/D_0)\xi,\qquad \xi=a(D)/a(D_0).
\]
Local square roots of $D/D_0$ at the conductor of $\eta$ give a value
$z$ of $\eta$ with $\eta(D/D_0)=z^2$.  Then $u(z)=z^e$ and
$\nu(D/D_0)=z^{e-1}$, so $\xi/z$ is fixed by every algebraic
automorphism and is rational.  Moreover the same local-square condition
gives $\psi(D/D_0)=\eta(D/D_0)=z^2$.  The comparison therefore reads
\[
 (\xi/z)^2=\frac{L_*(\sigma D)}{L_*(\sigma D_0)},
\]
which proves \eqref{co:exact-ratio} since $v(z)=0$.
There is no product over the new prime factors: the unramified
squarefree formula includes both local valuations zero and one.
No exceptional local representation at $2$ is excluded by the
same-local-class comparison.  The vanishing assertion and
\eqref{co:even-bound} follow.
\end{proof}

\subsection{A detecting Kummer functional}

Fix now the positive odd-sign filter and its original support, including
the auxiliary class-group tests needed for arithmetic realization.
Choose a fixed negative odd
fundamental discriminant $h_*$ with $L(E^{(h_*)},1)\neq0$, whose local
squareclasses agree with those of the filter at $2N$.  Adjoin its prime
factors to $S$, exclude them from the variable indices, and refine the
fixed unit tests there as required, all before varying $D$.
Choose it so that $|h_*|>4$ and one of its ramified primes
excludes the CM field of $E$, if there is one.  Thus
\[
 K_D=\mathbb Q(\sqrt{Dh_*})
\]
is imaginary, split at $2N$, and has only the units $\{1,-1\}$.
Fix a modular parametrization $\varphi$
with $\varphi(\infty)=0$ throughout.  With this normalization the
cuspidal constants in its Hecke and Atkin--Lehner relations are torsion;
see \cite{Manin1972,Drinfeld1973}.  An arbitrary translation by a
rational point is not permitted.

Let $H_D$ be the Hilbert class field of $K_D$.  Choose an
$\mathcal O_{K_D}$-ideal $\mathfrak n_D$ with
$\mathcal O_{K_D}/\mathfrak n_D\simeq\mathbb Z/N\mathbb Z$,
fixing one orientation at each level prime, and let
$x_D\in X_0(N)(H_D)$ be its conductor-one CM point.  The unramified
genus character $\chi_D$ attached to
$K_D(\sqrt{h_*})/K_D$ satisfies
$\chi_D([\mathfrak a])=\chi_{h_*}(\operatorname{Norm}\mathfrak a)$
on ideals prime to $Dh_*$.  Define the unaveraged character sum
\[
 P_D=\sum_{[\mathfrak a]\in\operatorname{Pic}(\mathcal O_{K_D})}
       \chi_D([\mathfrak a])\varphi(x_D^{[\mathfrak a]}).
\]
These conductor-one genus sums over the varying fields $K_D$ are
distinct from the conductor-$|h|$ ring-class sums over a fixed field
$K$ defined in Section~\ref{rl:section}.
The point $P_D$ lies in $E(\mathbb Q(\sqrt D,\sqrt{h_*}))$, with the quadratic
interpretation when $D=1$.  Its genus-character component over $K_D$
is the sum of the $D$- and $h_*$-twist components over $\mathbb Q$.
The latter has rank zero because $L(E^{(h_*)},1)\ne0$.
Thus the free part of $P_D$ lies in the $D$-twist component, a line
when $E^{(D)}$ has a simple zero.  We retain the unaveraged sum in the
theta comparison; a fixed integer $L$ kills bounded-degree torsion
and makes its comparison with a rational twist point integral.
At a simple zero write $j(P_D)$ for its index valuation in the free
rational twist line, allowing the fixed restriction and projection
indices.  Put $j(P_D)=+\infty$ if the product derivative vanishes.

\begin{lemma}[Kummer detection, including CM]\label{co:kummer}
Let $C$ act on $T=T_2E$ as complex conjugation, and choose a nonzero
integral row $b:T\longrightarrow\mathbb Z_2$ with $bC=b$.
For a fixed positive simple-zero witness $D_w$, there are arbitrarily
large precisions $M$ and good primes $\ell$ such that
\begin{equation}\label{co:ell-conditions}
 \rho_T(\operatorname{Fr}_\ell)\equiv C\pmod {2^M},\qquad
 (h_*/\ell)=-1,\qquad(D_w/\ell)=1.
\end{equation}
Evaluation at Frobenius squared defines an integral homomorphism
\[
 \lambda_M:E(\mathbb F_{\ell^2})\longrightarrow\mathbb Z/2^M\mathbb Z,
 \qquad \lambda_M=b\kappa_M,
\]
of Frobenius sign $+1$.  The valuation of
$\lambda_M(\operatorname{red}_\ell(LP_{D_w}))$ is bounded independently
of $M$, for a fixed nonzero integer $L$ making the twist comparison
integral.  Extra fixed finite-field restrictions can be retained when they agree
with the chosen complex-conjugation coset, or when a compatible lift
has been specified.
\end{lemma}

\begin{proof}
For every elliptic curve $A/\mathbb Q$, the rational Tate module $V_2A$
is irreducible.  Indeed Faltings' semisimplicity and endomorphism theorem
\cite[\S1, assertions (a),(b)]{Faltings1983} would turn a proper
summand into a nontrivial idempotent in
$\operatorname{End}_{\mathbb Q}(A)\otimes\mathbb Q_2=\mathbb Q_2$.
The equality of this endomorphism ring with $\mathbb Z$ includes CM
curves: a nonreal CM endomorphism is not defined over $\mathbb Q$.
The Tate image contains a nonidentity central scalar, by the homothety
theorem \cite[\S2, Theorem~3 and paragraph~7]{Bogomolov1981Points}.

Let $Q\in A(\mathbb Q)$ be nontorsion and let $c_Q$ be its division
cocycle.  On the kernel of the Tate action its image is a closed
$\mathbb Z_2$-submodule of $T_2A$, stable under the full Galois image.
If its rational span were zero, the rational Kummer class would factor
through the Tate image.  A central scalar $z\neq1$ annihilates its
first cohomology: for a cocycle $c$ the equality of $c(zg)$ and $c(gz)$
gives $(z-1)c(g)=(g-1)c(z)$, so $c$ is a coboundary over $\mathbb Q_2$.
This contradicts Kummer injectivity and the finite generation of
$A(\mathbb Q)$.  Irreducibility consequently makes the translation
image contain $2^aT_2A$ for some $a$.  Intersecting with the kernel of
additional finite data retains a full lattice.

Apply this to the rational twist point supplied by $LP_{D_w}$ on
$A=E^{(D_w)}$.  Since $D_w>0$, conjugation acts on its Tate module as
$C$.  Modify a conjugation lift by an element of the simultaneous
kernel whose translation is $v$.  The square of the modified element
has translation $(1+C)v$, and $b(1+C)v=2b(v)$.  The full-lattice
property allows a choice with a fixed finite nonzero valuation.
Conjugation acts trivially on $\mathbb Q(\sqrt{D_w})$ and nontrivially
on $\mathbb Q(\sqrt{h_*})$.  Chebotarev in each finite torsion,
division, and character extension gives \eqref{co:ell-conditions}
and the chosen finite nonzero detection.

For good $\ell$, Lang's theorem and finite-field Kummer theory give
\[
 E(\mathbb F_{\ell^2})/2^M E(\mathbb F_{\ell^2})
 \simeq T/\bigl(2^MT+(\operatorname{Fr}_\ell^2-1)T\bigr)=T/2^MT.
\]
This is evaluation of a cocycle at arithmetic Frobenius squared.
Its equivariance and $bC=b$ prove the sign assertion.  Reduction of
the division cocycle gives precisely the detection already arranged.
Finally $\det C=-1$ implies $\ell\equiv-1\pmod {2^M}$.
\end{proof}

\subsection{The finite Schwartz function and its transformation}

\begin{lemma}[Explicit weighted theta construction]\label{co:theta-bridge}
For $\ell$ as in Lemma~\ref{co:kummer}, the reduced modular
parametrization followed by $\lambda_M$ admits an integral weighted
ternary theta lift.  After extraction of the coefficients at multiples
of $|h_*|$ it gives a classical holomorphic weight-$3/2$ form
$g_M=\sum a_M(n)q^n$, of level supported on $S\cup\{\ell\}$, with
\eqref{co:recurrence} modulo $2^M$.  The associated characters have
$2$-power order and quadratic ratio.  For squarefree admissible $D$,
\begin{equation}\label{co:weighted-comparison}
 a_M(D)=c_0\,2^{\omega(N)}\zeta_M(D)
       \lambda_M\bigl(\operatorname{red}_\ell(LP_D)\bigr)
                     \pmod {2^M},
\end{equation}
where $\zeta_M(D)$ is a $2$-power root of unity and $c_0,L$ are fixed
integers independent of $D,M,\ell$.  The extra index restriction at
$\ell$ is $(D/\ell)=1$.
\end{lemma}

\begin{proof}
Let $B$ be ramified at $\ell$ and infinity, $R$ an Eichler order of
level $N$, and $V=B^0$ with $Q(x)=-x^2=\operatorname{Nrd}(x)$.
The root characters at the fixed level primes below are chosen once.
Let $\mathcal O$ be the valuation ring at a place above $2$ of the
field generated by their values.  The other local weights are
quadratic, so this field is independent of $M$ and $\ell$.
In this construction integrality means membership in $\mathcal O$,
and congruences mean congruences modulo $2^M\mathcal O$.
\par\smallskip\noindent\emph{The supersingular coefficient function.}
The supersingular pairs on $X_0(N)$ are the finite set
\[
 X=B^\times\backslash\widehat B^\times/\widehat R^\times.
\]
This description permits every $N$ prime to $\ell$, including arbitrary
depth at $2$; see \cite[\S\S1.1--1.2,2.4--2.6]{JetchevKane2010}.
For $z\in X$ put $R_z=B\cap z\widehat Rz^{-1}$.
Choose $L$ killing the finite subgroup generated by the images of the
cuspidal constants in the Atkin--Lehner relations, and also all the
bounded-degree torsion used below.  Then
\[
 P_M(z)=\lambda_M(L\varphi_\ell(z)),
\]
where $\varphi_\ell$ is the specialization of $\varphi$ to the
good-reduction special fiber.  This function has Brandt eigenvalues
$a_p(E)$ away from $N\ell$, Atkin--Lehner
eigenvalue $w_{p^e}$ at $p^e\Vert N$, and eigenvalue $+1$ under the
right $\ell$-uniformizer.  These identities are modulo $2^M$.
Choose any integer lifts of this finite list of values.

\par\smallskip\noindent\emph{The local Schwartz weights.}
We now define the local functions, rather than appealing to an
unspecified weighted theta theorem.  Off $2Nh_*\ell$ take the
characteristic function of $R_p^0$.  At $p^e\Vert N$, use a model
\[
 R_p=\left\{\begin{pmatrix}a&b\\c&d\end{pmatrix}
       :a,b,d\in\mathbb Z_p, c\in p^e\mathbb Z_p\right\}.
\]
Require the split fundamental unit squareclass and an integral
embedding.  Its action on the cyclic level quotient selects one root
$r\in\mathbb Z_p^\times$ of $x^2$.  Give $x$ weight $\alpha_p(r)$,
where $\alpha_p$ is a finite unit character of $2$-power order with
\begin{equation}\label{co:root-parity}
 \alpha_p(-1)=w_{p^e}\chi_{h_*}(p)^e.
\end{equation}
Such a character exists for either prescribed sign.  The root is a
locally constant function to any fixed character precision: the two
idempotents of the split integral quadratic algebra distinguish the
two summands of the cyclic quotient.  Unit scaling $x\mapsto ux$
multiplies this weight by $\alpha_p(u)$, and order-unit conjugation
preserves it.

At $2$ require $(1+x)/2\in R_2$ and $x^2\in\mathbb Z_2^{\times2}$.
If $2\mid N$, the preceding selected-root rule is imposed using
$e_x=(1+x)/2$.  Its two eigenvalues $(1+r)/2,(1-r)/2$ differ by the
odd unit $r$; hence the two idempotents remain distinct modulo $2$.
This verifies the construction at every exponent of $2$ in $N$.
Scaling by an odd $u$ preserves integrality, since
$(1+ux)/2=u(1+x)/2+(1-u)/2$.  This is an invertible affine change of
the operator $e_x$, so it preserves its two summands individually.
The selected summand is selected by the cyclic level quotient, not
by declaring its eigenvalue to be $0$ or $1$ modulo $2$.  On that same
summand $x$ acts as $r$ and $ux$ acts as $ur$.  Even when
$u\equiv-1\pmod4$ interchanges the two eigenvalue parities, the root
weight therefore changes by exactly $\alpha_2(u)$.
If $2\nmid N$ there is no root weight at $2$.

For $p\mid h_*$ require $x\in M_2(\mathbb Z_p)^0$ and
$v_p(Q(x))=1$.  Its reduction is a nonzero nilpotent.  Fix
$X_0=\left(\begin{smallmatrix}0&1\\0&0\end{smallmatrix}\right)$ and
define
\[
 \epsilon_p(gX_0g^{-1})=\left(\frac{\det g}{p}\right).
\]
The centralizer consists of
$\left(\begin{smallmatrix}a&b\\0&a\end{smallmatrix}\right)$, with
determinant $a^2$, so this is well-defined on the entire nonzero
nilpotent orbit.  It obeys
\[
 \epsilon_p(uyu^{-1})=(\det u/p)\epsilon_p(y),\qquad
 \epsilon_p(ay)=(a/p)\epsilon_p(y).
\]
The second equality follows by conjugating $X_0$ by
$\operatorname{diag}(a,1)$.  Set $f_p(x)=\epsilon_p(\bar x)$ on the
specified support.  Thus the right conjugation character is precisely
the local component of $\chi_{h_*}\circ\operatorname{Nrd}$.

At $\ell$, let $\mathfrak P$ be the maximal ideal of the maximal
order.  Its residue field is $\mathbb F_{\ell^2}$, and the reduction of
an integral trace-zero element lies on its one-dimensional trace-zero
line.  Fix a nonzero basis vector $\beta$ on that line and put
\[
 f_\ell(x)=\left(\frac{t}{\ell}\right)
   \quad\text{if }\bar x=t\beta\neq0,
 \qquad f_\ell(x)=0\quad\text{otherwise}.
\]
Order units act trivially by conjugation on the residue field.
A uniformizer acts by its nontrivial automorphism and sends
$\beta$ to $-\beta$.  Therefore its multiplier is $-1$, since
$\ell\equiv-1\pmod4$; scaling has character $(u/\ell)$.
Moreover $\beta^2$ is a nonsquare and $-1$ is a nonsquare modulo
$\ell$, so $Q(x)=-t^2\beta^2$ is a nonzero square.  Because
$(|h_*|/\ell)=1$, the extracted indices automatically satisfy
$(D/\ell)=1$.  No average over $\mathbb F_\ell^\times$ is used.

All these functions have compact support and are locally constant.
Let $f=\prod f_p$, incorporating any additional fixed unit-squareclass
conditions by their characteristic functions.  Those conditions are
stable under unit squares and change none of the scaling characters.
For $u\in\widehat R^\times$,
\begin{equation}\label{co:order-equivariance}
 f(u^{-1}xu)=\chi_{h_*}(\operatorname{Nrd}u)^{-1}f(x).
\end{equation}
The product scaling character is even.  Indeed its value at $-1$ is
\[
 \left(\prod_{p^e\Vert N}w_{p^e}\chi_{h_*}(p)^e\right)
       \chi_{h_*}(-1)(-1)=1;
\]
the parenthesized factor is the even functional sign of $E^{(h_*)}$.
Denote its associated classical character by $\eta_B$.  In the
inverse-idele convention fixed above, if $\alpha$ is this finite
scaling character, then $\eta_B(d)=\alpha(d)$ for positive integers
$d$ prime to the support: the inverse idele character becomes the
product of the finite unit values at the lower-right entry.

\par\smallskip\noindent\emph{The scalar transformation law.}
Here are all the Weil-representation checks needed for scalar
modularity.  For an additive character of conductor $\mathbb Z_p$,
write $n(b)=\left(\begin{smallmatrix}1&b\\0&1\end{smallmatrix}\right)$,
$m(u)=\operatorname{diag}(u,u^{-1})$, and
$n^-(c)=\left(\begin{smallmatrix}1&0\\c&1\end{smallmatrix}\right)$.
Use a Weyl lift $w$, its compatible self-dual Fourier transform, and
the Weil factors $\gamma_p(u),\gamma_p$ for this additive character.
The defining formulas are
\[
 \omega(n(b))f_p(x)=\mathbf e_p(bQ(x))f_p(x),\quad
 \omega(m(u))f_p(x)=\gamma_p(u)|u|_p^{3/2}f_p(ux),\quad
 \omega(w)f_p=\gamma_p\widehat f_p.
\]
They follow from the Weil construction and Poisson summation
\cite[I, no.~13]{Weil1964}.  Integral $b$ fixes $f_p$, since $Q$ is integral on
its support.  If $f_p$ is invariant under translations by a lattice
$L_p'$, its Fourier transform is supported on the dual lattice
$(L_p')^\vee$.  Choose $A_p$ so that
$p^{A_p}Q((L_p')^\vee)\subset\mathbb Z_p$.  Then $n(p^{A_p}b)$ fixes
$\widehat f_p$, and conjugating by $w$ proves that every
$n^-(p^{A_p}b)$ fixes $f_p$.  This argument works without change at
$2$ and for all finite character depths.

Relative to the ordinary theta multiplier cubed, the diagonal Weil
factor is $(\det Q,u)_p$.  Indeed diagonalizing $Q$ and multiplying
the one-dimensional Weil-index identities
\cite[II, nos.~25, 28]{Weil1964} gives the Hilbert symbol
of the product of its three coefficients.  In a presentation
$B=(-a,-b)$ that product is the determinant of
$aX^2+bY^2+abZ^2$, namely $(ab)^2$.  Thus this additional factor is
$1$.  The remaining diagonal action is exactly the scaling character
already computed.  Increasing $A_2$ to the depth needed for the
standard metaplectic splitting, the factorization
\[
 \begin{pmatrix}a&b\\c&d\end{pmatrix}
    =n^-(c/a)m(a)n(b/a),\qquad c\in p^{A_p}\mathbb Z_p, a\in\mathbb Z_p^\times,
\]
proves scalar transformation on the required local $\Gamma_0$.
Conjugation by $z$ commutes with these Weil actions and preserves $Q$.
Consequently the same level and scalar character work for every
summand.  The Gaussian at infinity gives holomorphic weight $3/2$;
positive definiteness gives nonnegative exponents at every cusp.

Choose an integer $c_0$ divisible by every possible $|R_z^\times|$.
Such a uniform choice exists because these unit groups embed in
$\operatorname{GL}_4(\mathbb Z)$ and their orders are bounded.
The theta series is
\begin{equation}\label{co:theta-sum}
 H_M(n)=c_0\sum_{z\in X}
  \frac{P_M(z)\chi_{h_*}(\operatorname{Nrd}z)}{|R_z^\times|}
  \sum_{\substack{x\in B^0\\Q(x)=n}} f(z^{-1}xz).
\end{equation}
Equation~\eqref{co:order-equivariance} proves independence of right
representatives; rational reduced norms are positive, so the norm
character is trivial on the left rational action.  The preceding Weil
calculation proves modularity \cite[III, Theorem~6]{Weil1964}, and
the choice of $c_0$ proves integral Fourier coefficients.  More
precisely, each $c_0/|R_z^\times|$ is an integer, so changing a lift
of $P_M(z)$ by $2^M$ changes every coefficient by an element of
$2^M\mathcal O$.  No precision is lost to a stabilizer denominator.
The chosen integer lifts need satisfy no characteristic-zero
eigenrelation.

\par\smallskip\noindent\emph{The good-prime Hecke relation.}
For completeness, the good-prime relation can be checked in a matrix
order.  Neighbors correspond to the lines in $\mathbb F_p^2$.
A nonzero primitive nilpotent reduction has exactly one invariant
line.  For an integral element of nonzero reduction and norm $n$,
the number of invariant lines is $1+(-n/p)$; zero reduction contributes
the additional $p$ lines.  If a primitive $y$ has $p^2\mid Q(y)$,
choose a basis in which its reduction is $X_0$.  Its lower-left entry
is divisible by $p^2$, and exactly the neighbor for that invariant
line makes $y/p$ integral.  These counts partition the three terms
of the theta Hecke operator.  Scaling by $p$ at the support contributes
$\eta_B(p)$, and the reduced norm of a Brandt neighbor contributes
$\chi_{h_*}(p)$.  Thus, with $A(n)=H_M(n)$,
\[
 A(p^2n)+\eta_B(p)(-n/p)A(n)+p\eta_B(p)^2A(n/p^2)
       =\eta_B(p)\chi_{h_*}(p)a_p(E)A(n)\pmod {2^M}.
\]
The orbit-stabilizer weights in \eqref{co:theta-sum} are exactly the
weights in this neighbor count.  Extracting
$a_M(n)=H_M(|h_*|n)$ gives
$\eta=\eta_B\chi_{h_*}$ and
$\psi=\eta_B(|h_*|/\cdot)$, proving \eqref{co:recurrence}.
Fixed-support coefficient sieves are exact zero-or-retain operations;
they preserve coefficient integrality and commute with this relation.

\par\smallskip\noindent\emph{The weighted coefficient comparison.}
It remains to prove the weighted comparison; modularity alone does
not supply it.  At $n=|h_*|D$ the counted elements are precisely the
optimal embeddings $\mathcal O_{K_D}\hookrightarrow R_z$, sending
$\sqrt{Dh_*}$ to $x$.  The condition at $2$ is exactly the condition
on the integral generator $(1+\sqrt{Dh_*})/2$.
At a level prime there are two orientations.  Indeed a stable lattice
splits into its two integral idempotent summands, and a cyclic quotient
of order $p^e$ lies wholly in one summand.  At $\ell$ there are two
orientations, the two residue-field embeddings.  At all remaining
primes there is a single local optimal-embedding orbit.

Fix a rational embedding $x_0$ and an initial orientation tuple $z_0$.
The orbits with these orientations are indexed by
$\operatorname{Pic}(\mathcal O_{K_D})$, with representatives $az_0$
for $a\in\widehat K_D^\times$.  The equality
\[
 (az_0)^{-1}x_0(az_0)=z_0^{-1}x_0z_0
\]
shows that the Schwartz weight is constant along that orbit.  Its
remaining relative weight is
$\chi_{h_*}(\operatorname{Norm}_{K_D/\mathbb Q}a)$, exactly the
unramified genus character.  The reduction correspondence commutes
with this ideal-class action and remembers the cyclic level subgroup
\cite[\S2]{JetchevKane2010}; a Galois-equivariant description of the
lift, class action and reduction is given in
\cite[\S3.2, \S\S4.6, 4.9--4.10 and Corollary~4.1]{CornutJetchev2013}.
We use these correspondences only for their geometry; the weight
identity has just been computed, and the multiplicity is computed next.

An embedding stabilizer is $\mathcal O_{K_D}^\times=\{\pm1\}$.
An orbit in $R_z$ therefore contains $|R_z^\times|/2$ elements, giving
the factor $1/2$ in \eqref{co:theta-sum}.  At $p^e\Vert N$, reversing
orientation multiplies its Schwartz factor by
$\alpha_p(-1)=w_{p^e}\chi_{h_*}(p)^e$, and its norm-weighted
parametrization value by the same sign.  At $\ell$ both multipliers
are $-1$: one is the residue-line calculation and the other is
$\chi_{h_*}(\ell)$ times the Frobenius eigenvalue $+1$.
Every one of the $2^{\omega(N)+1}$ orientation tuples consequently
has the same contribution.  The initial tuple has a $2$-power unit
phase.  Combining this count with its factor $1/2$ gives exactly
\eqref{co:weighted-comparison}.  There is no division by a class
number, and none of its rational multipliers grows with $\ell$ or
the number of factors of $D$.
\end{proof}

\begin{proposition}[Odd coefficient detection]\label{co:odd}
The forms in Lemma~\ref{co:theta-bridge}, after a fixed adjustment of
precision, satisfy
\[
 v(a_M(D))\geq\min\{M,j(P_D)-C\},\qquad
 v(a_M(D))<M\ \Longrightarrow\ j(P_D)\leq v(a_M(D))+C.
\]
A visible coefficient detects $a(E^{(D)})=1$, and a fixed positive
witness has bounded coefficient valuation as $M\to\infty$.
All constants are independent of $M,\ell,D$.
\end{proposition}

\begin{proof}
Gross--Zagier identifies the height of the genus sum with the product
derivative \cite{GrossZagier1986}.  Since $L(E^{(h_*)},1)\neq0$,
vanishing of the derivative makes $P_D$ torsion, and a nontorsion
$P_D$ is equivalent to a simple zero of $E^{(D)}$.
The genus sum lies in a biquadratic field, so a fixed integer kills
all possible torsion: reduce prime-to-residue-characteristic torsion
at two fixed good primes of different residue characteristic; the
residue degrees are at most four.  Include that integer in $L$.
On the free twist line, divisibility by $2^{j(P_D)}$ remains divisibility
after the integral maps of reduction and $\lambda_M$.
Restriction and the two quadratic projections have composite
multiplication by $2$, so their integral indices have a fixed bound.
Equation~\eqref{co:weighted-comparison} therefore proves both
valuation assertions, with fixed losses only.  Lemma~\ref{co:kummer}
gives the witness.
\end{proof}

\subsection{Minimum depth and integral coefficient operators}

We now pass from detection to normalization.  For the even source,
the cyclotomic lower bound gives \eqref{co:even-bound} at every
nonzero squarefree coefficient.  For the odd source, assume for now
that every simple-zero parameter satisfies
\[
 j(P_D)\geq w(D)+s(D)/2-C_*.
\]
Proposition~\ref{pw:odd-lower} will prove this using only the even
construction.  Proposition~\ref{co:odd} then gives
\[
 v(a_M(D))\geq\min\{M,w(D)+s(D)/2-C'\}
\]
at simple zeros, where $s(D)$ is defined.  At every other odd-sign
parameter the product derivative vanishes and $a_M(D)=0$ modulo
$2^M$.  In particular the coarser bound
$v(a_M(D))\geq\min\{M,w(D)-C'\}$ holds at every admissible $D$.
The remaining coefficient calculus is conditional on these lower
bounds; it does not establish them.

For the operator and graph arguments, at each actual odd stage we
write $S$ for the original fixed support together with the primes of
$h_*$ and the current reduction prime $\ell$.  The latter is stationary
during that stage and is never an element of a newly adjoined set
$U$.  Thus every statement below about ramification outside $S\cup U$
includes $\ell$.  This notation does not change the original fixed
data controlling the constants, or choose a new minimum when a vertex
is omitted.

\begin{lemma}[Stabilized minimum depth]\label{co:minimum}
Assume the uniform arithmetic lower bounds just stated.  Each coefficient source
defines a binary symbol
\begin{equation}\label{co:symbol}
 A(D)=\begin{cases}
 1,&v(a(D))=\delta+w(D),\\
 0,&v(a(D))>\delta+w(D)\text{ or }D\text{ is off the filter},
 \end{cases}
\end{equation}
with a fixed minimum normalized depth $\delta$ and at least one
bounded-weight witness.  For the odd source $a=a_M$ at an actual stage:
the numerical minimum and its divisibility statements stabilize at
every fixed finite weight bound.  The finite symbol tables on arithmetic
labels will be stabilized separately in Section~\ref{pw:odd-stages}.
At a unit symbol,
\[
 v(L_*(\sigma D))\leq2w(D)+O(1)
 \quad\text{or}\quad j(P_D)\leq w(D)+O(1),
\]
respectively, and the analytic order is respectively zero or one.
\end{lemma}

\begin{proof}
For the even source normalize $a(D_0)=1$; the phase-removed squarefree
coefficients are rational by Proposition~\ref{co:even}.  For the odd
source, \eqref{co:weighted-comparison} makes every visible squarefree
coefficient an integer times a root of unity modulo $2^M$.
Their finite valuations are therefore integers in both cases, not
arbitrary elements of a ramified coefficient field's value group.
The fractional part of $w(D)$ is fixed on the filter,
because the parity of the number of simple primes is its fixed
quadratic discriminant character.  The lower bounds and a nonzero
witness therefore give a minimum for the even source.

For the odd source choose a nonprincipal ultrafilter on the increasing
precisions.  Minimize $v(a_M(D_M))-w(D_M)$ over sequences with bounded
weight and visible coefficient on an ultrafilter-large set.  The
uniform lower bound, discrete values, and fixed witness give a finite
attained minimum.  At each fixed weight cutoff $W$, failure of its
divisibility on an ultrafilter-large set would allow a choice of a
violating $D_M$ with $w(D_M)\leq W$, contrary to the minimum.
Hence every prescribed finite list of these divisibilities holds at
a sufficiently advanced actual stage.  If larger weights arise in a
coefficient-integrality check, the coarse bound
$\min\{M,w(D)-C\}$ already pays any fixed normalization once both
$M$ and $w(D)$ exceed its fixed threshold.  This also explains why
no infinite-precision modular lift is asserted or needed.
The last conclusions follow from the two coefficient comparisons.
\end{proof}

Write $U_p\sum c(n)q^n=\sum c(pn)q^n$ and
$V_p\sum c(n)q^n=\sum c(n)q^{pn}$.  For a squarefree set $U$ put
\begin{equation}\label{co:BU}
 B_U=2^{-\lceil\delta+w(U)\rceil}
       \prod_{p\in U}(U_p-\eta(p)V_p)g,\qquad
 e_U=\lceil\delta+w(U)\rceil-\delta-w(U).
\end{equation}
We call $U$ matching if $e_U=0$ and opposite if $e_U=1/2$.
These operators introduce each new prime to level exponent one:
$V_p$ changes the character by $(p/\cdot)$ and raises the level by
$p$, while $U_p=U_{p^2}V_p$ on that raised level.
Thus the new character is $\psi_U=\psi\prod_{p\in U}(p/\cdot)$.

\begin{lemma}[Coefficient bounds and unary removal]\label{co:operators}
For the even source let $U$ be arbitrary.  For the odd source, fix
an integer $R\ge1$ and a weight bound $W$; all assertions below are
made at a sufficiently advanced actual stage, simultaneously for every
$U$ with $w(U)\le W$ and every Fourier index.
In these respective settings the coefficients of $B_U$ are integral.
Write uniquely
$n=Jrm^2$, with $J\subseteq U$, $r$ squarefree, and $(r,U)=1$.
The coefficient vanishes unless $D=(U/J)r$ belongs to the filter,
and for the even source otherwise
\begin{equation}\label{co:BU-bound}
 v(B_U[n])\geq w(r)+\tfrac12\#\{p\in J:p\text{ simple}\}-e_U.
\end{equation}
At that stage the odd source satisfies this inequality for every
$U$ with $w(U)\le W$ and every $n$, with its right side replaced
by its minimum with $R$.
In particular the untruncated bound holds on any prescribed bounded
range of squarefree weights, after choosing $R$ and the stage.
For a split $p\in J$ the difference factor, after its unit phase and
one factor of $2$ are removed, has residue
\begin{equation}\label{co:cp}
 c_p(D)=\frac{a_p(E)-(\psi/\eta)(p)(-D/p)-1}{2}\pmod2
\end{equation}
at every square shell.  For a simple $p$ this formula holds on the
first shell $p\nmid m$.

Choose $i\in\{0,1\}$ such that $\eta\chi_{-4}^i$ is odd, and set
\[
 \Theta_J=\sum_{\substack{m>0\\(m,S)=1}}
              \eta(m)\chi_{-4}(m)^i m q^{Jm^2}.
\]
For matching $U$ the modular form
\begin{equation}\label{co:HU}
 H_U=\frac12\left(B_U-
       \sum_{\substack{J\subseteq U\\J\text{ split only}}}
                 B_U[J]\Theta_J\right)
\end{equation}
has integral coefficients, with the same stage and weight restriction
in the odd case.  For opposite $U$ put $H_U=B_U$.
The new primes in the level of $\Theta_J$ have exponent one, and its
character is $\eta\chi_{-4}^i(-J/\cdot)$.
\end{lemma}

\begin{proof}
After removing the powers of $\eta(p)$, the square-index multipliers
in \eqref{co:recurrence} are
\[
 b_0=1,\quad b_1=a_p-\kappa,\quad
 b_{s+1}=a_pb_s-pb_{s-1},\qquad
 \kappa=(\psi/\eta)(p)(-D/p).
\]
If $p\mid D$, then $\kappa=0$ and $b_s$ alternates between $1$ and
$0$ modulo $2$, because $a_p$ is even.  Thus the difference from
$p\in U/J$ is odd.  If $p\nmid D$, all $b_s$ are odd and the
difference from $p\in J$ is even.  Subtracting the normalization in
\eqref{co:BU} gives exactly
\[
 \delta+w((U/J)r)+\#J-\lceil\delta+w(U)\rceil
 =w(r)+\tfrac12\#\{p\in J:p\text{ simple}\}-e_U.
\]
This is an integer on the allowed filter.  Its only potential
negative value before imposing parity is $-1/2$, so it is nonnegative.

For the odd source this calculation first applies on a fixed bounded
range of squarefree weights, with enough precision to pay the fixed
normalizing power of $2$.  To obtain the global truncated assertion,
fix $R$ and a bound $W$ for $w(U)$.  The normalizing exponents
$\lceil\delta+w(U)\rceil$ are uniformly bounded.  Above a sufficiently large weight
threshold, the coarse bound $\min\{M,w(D)-C\}$ from
Lemma~\ref{co:minimum} pays that normalizing power and at least $R$
further factors of $2$, once $M$ is large enough.  Below the threshold
use the stabilized minimum and increase the precision again if needed.
The recurrence has integral multipliers, so errors modulo $2^M$
remain divisible by $2^M$ at every square shell before normalization.
Thus at the chosen stage the truncated bound holds at every
coefficient for every $U$ with $w(U)\le W$, giving integrality
without asserting an unbounded valuation estimate for a fixed
finite-precision lift.

For split $p$, rationality of $E[2]$ over $\mathbb F_p$ gives
$a_p\equiv p+1\pmod4$.  With $\Delta_s=b_{s+1}-b_s$, every
$\Delta_s$ is even and
$\Delta_{s+1}=a_p\Delta_s-p\Delta_{s-1}$, so
$\Delta_{s+1}\equiv\Delta_{s-1}\pmod4$.
Direct calculation gives
\[
 \Delta_1-\Delta_0
  =a_p^2-a_p\kappa-2a_p+2\kappa+1-p\equiv0\pmod4.
\]
Consequently every $\Delta_s/2$ has the residue in \eqref{co:cp}.
At a simple prime only $\Delta_0$ is used.

In matching parity the right side of \eqref{co:BU-bound} can vanish
only when $r=1$ and $J$ is split only.  The same recurrence, retaining
the character phases, gives
$B_U[Jm^2]\equiv B_U[J]\eta(m)\pmod {2}$.
The corresponding coefficient of $\Theta_J$ has exactly that
residue, and different squarefree $J$ have disjoint supports.
The numerator in \eqref{co:HU} is therefore divisible by the integer
$2$, not just by the maximal ideal of a ramified character field.
The ordinary odd-character unary theta transformation and dilation
by $J$ give the asserted character and level.  For the odd source
these calculations use the same sufficiently advanced stage,
uniformly for $w(U)\le W$, with precision paying the fixed
divisions as in Lemma~\ref{co:minimum}.
\end{proof}

\subsection{Trace operators and omission of a whole rational prime}

The next step turns changing a Frobenius label into a difference of
coefficient tests.  Fix a coefficient valuation ring $\mathcal O$,
enlarging its field if needed for the finite eigensystem span below,
and write $\mathfrak m$ for its maximal ideal.  Divisibility by $2^a$
means divisibility in $2^a\mathcal O$; reduced traces and expansions
mean coefficientwise reduction of Fourier expansions modulo $\mathfrak m$.
These are different operations; no saturation of the integral module
below is assumed.
The normalized squarefree tests are binary because their phase-removed
values are rational $2$-adic integers (integer congruences at an odd
stage), and every $2$-power root of unity reduces to $1$.

Multiply the form under consideration by
$\theta=1+2\sum_{m>0}q^{m^2}$.  This gives a weight-two form with the
same reduced expansion.  Work in the finite-dimensional sum of its
characteristic-zero good Hecke eigensystems, retaining old
multiplicity spaces and its diamond transforms.  The associated
two-dimensional representations are the usual modular Galois
representations for cusp forms \cite{Deligne1969}, and the sums of
the two defining characters for Eisenstein series.  Their commuting
trace and determinant operators are denoted by $t(u),d(u)$.
Arithmetic Frobenius at a good prime $q$ has trace $T_q$ and
determinant $q\langle q\rangle$.
Let $\mathcal M$ be the $\mathcal O$-module generated by this
weight-two vector and the undivided terms $\theta B_U[J]\Theta_J$,
together with all their good-Hecke and diamond translates.

\begin{lemma}[Integral traces]\label{co:traces}
The module $\mathcal M$ is an integral Fourier lattice preserved by
the trace and determinant operators.  On the coefficientwise reduced
trace orbit of the indicated vector,
\begin{equation}\label{co:trace-identities}
 d(u)t(v)=t(v),\qquad
 t(u)t(v)=t(uv)+t(uv^{-1}),\qquad t(1)=0.
\end{equation}
Every coefficient test of a finite product of traces can be performed
using distinct fresh primes avoiding the tested index.  In particular
testing $t(u_1)\cdots t(u_a)$ at $n$ becomes testing the original
reduced coefficient at $nq_1\cdots q_a$.
\end{lemma}

\begin{proof}
We first check integrality of the generators.  Diamonds act by scalars
on $B_U$ and on each unary term.  For \eqref{co:HU}, the only apparent
denominator is removed by
\[
 \langle u\rangle H_U=\psi_U(u)H_U+
 \frac12\sum_J B_U[J]
 \bigl(\psi_U(u)-\eta(u)\chi_{-4}(u)^i(-J/u)\bigr)\Theta_J;
\]
the difference of the two characters is zero or twice a unit.
The same formula applies after multiplication by theta.
Good Hecke operators commute with diamonds, and their coefficient
formula preserves integrality.  Thus all generators of $\mathcal M$
are integral.  Finitely many Fourier coefficients inject the ambient
finite-dimensional space into a finite-dimensional coefficient space,
so $\mathcal M$ embeds in $\mathcal O^r$ for some finite $r$.
It is consequently a finite, closed $\mathcal O$-module.
Continuity and simultaneous Chebotarev approximation on the finitely
many eigensystems approximate $t(u)$ and $d(u)$ by $T_q$ and
$q\langle q\rangle$ at every prescribed precision.  These operators
preserve $\mathcal M$, so their limits do as well.  They act by scalars
on each good eigensystem and by the identity on its old multiplicity
space, which retains the asserted commutation with level operators.

The formula
$(T_qf)[n]=f[qn]+q(\langle q\rangle f)[n/q]$ has no second term
when $q\nmid n$.  Successive distinct fresh primes therefore give
the coefficient assertion.  On a single character system,
determinants reduce to $1$ because all characters have $2$-power
order and the cyclotomic character is odd.  On a subtracted vector
the displayed diamond formula leaves unary corrections.  A fresh
trace kills each correction modulo the maximal ideal: the tested
index $nq$ cannot have its prescribed squarefree part $J$ when
$q\nmid nJ$.  This continues to hold after further fresh tests.
Thus $d(u)t(v)=t(v)$ on the reduced trace orbit.  Apply the matrix
identity
\[
 \operatorname{tr}(A)\operatorname{tr}(B)
  =\operatorname{tr}(AB)+\det(B)\operatorname{tr}(AB^{-1})
\]
to each characteristic-zero system and use that determinant
identity.  Finally $t(1)=2$ before reduction.
\end{proof}

\begin{lemma}[The level-prime inertia identity]\label{co:inertia}
For $p\in U$, choose inertia $\tau_p$ on which the new ramified
quadratic character is $-1$, and an arithmetic Frobenius lift $\phi_p$
in the same decomposition group.  In characteristic zero,
\begin{equation}\label{co:inertia-identity}
 U_p^{(2)}\bigl(t(g\tau_p\phi_p)-t(g\phi_p)\bigr)
   =d(\phi_p)\bigl(t(g\tau_p)-t(g)\bigr).
\end{equation}
Here $U_p^{(2)}$ extracts $p$-multiples in weight two.  The identity
commutes with further good traces.
\end{lemma}

\begin{proof}
A constituent old from a level omitting $p$ has trivial inertia, so
both differences are zero, whatever $U_p^{(2)}$ does on its old
multiplicity space.  Every other constituent has conductor exponent
one and primitive ramified quadratic determinant.  The conductor
formula excludes monodromy and identifies it as principal series
with one unramified and one ramified character.  Local compatibility
\cite[Theorem~(A)]{Carayol1986}, after converting its geometric
Frobenius convention to the arithmetic convention fixed above,
identifies its $U_p^{(2)}$ eigenvalue
with the Frobenius eigenvalue $\lambda$ on the unramified line.
Write
\[
 \rho(\tau_p)=\begin{pmatrix}1&0\\0&-1\end{pmatrix},\qquad
 \rho(\phi_p)=\begin{pmatrix}\lambda&0\\0&\mu\end{pmatrix}.
\]
If $z$ is the lower-right entry of $\rho(g)$, the two trace
differences are $-2z\mu$ and $-2z$.  Both sides of
\eqref{co:inertia-identity} equal $-2z\lambda\mu$.
The Eisenstein character pair has the same calculation.
A unary term is old if $p\notin J$, and has precisely this ramified
quadratic character if $p\in J$.  This exhausts the constituents.
\end{proof}

For the interface below a \emph{label} is the restriction of a
Frobenius element to an elementary $2$-extension unramified outside
the indicated support.  Its coordinates are quadratic characters.
When $F$ is quadratic, choosing one of the two places at a split prime specifies an
orientation; changing its ramification bit means multiplying by the
image of inertia at that place.  Product labels are the restrictions
of products of Frobenius elements, so their quadratic coordinates add.
Simple primes use rational quadratic labels.  Section~\ref{gr:section}
will realize these labels simultaneously by primes and interpret the
independent inertia changes as graph edges.

\begin{proposition}[Symbol interface]\label{co:symbol-interface}
Under the arithmetic lower bounds in Lemma~\ref{co:minimum}, the
symbols of either coefficient source have the following properties.
\begin{enumerate}
\item For matching $U$, the trace on split elements factors through
the maximal elementary $2$-extension of $F$ unramified outside
$S\cup U$.  For opposite $U$, traces factor through the corresponding
rational elementary extension.  Adding a split vertex $q$ is read
by coefficient $1$ of $t(\operatorname{Fr}_q)H_U$; adding a simple
vertex is read on the opposite $B_U$.
\item Let $p\in U$ and $g=\operatorname{Fr}_q$.  Changing one
oriented ramification bit at $p$ changes the symbol by
\begin{equation}\label{co:contraction}
 A\bigl((U\setminus\{p\})\cup\{q'\}\bigr),\qquad
 \operatorname{Fr}_{q'}\sim g\phi_p,
\end{equation}
unless $p,q$ are both simple, in which case the change is zero.
The replacement has the product labels at all retained places.
Its coefficient system has support $S\cup(U\setminus\{p\})$:
the entire rational prime $p$, including both places above it when
it splits in $F$, is omitted.
\item A split vertex with trivial label in the elementary extension
for the other vertices has zero symbol.
\item In matching parity, insertion of two simple primes with equal
rational labels on the fixed support and at the existing split
vertices preserves the symbol if
$\#E(\mathbb F_q)/2$ is odd for that label.  Their mutual simple
bits and their incidences with other simple vertices do not affect
this statement.
\end{enumerate}
These are coefficient identities.  Simultaneous realization of
their labels by primes is a separate arithmetic construction.
\end{proposition}

\begin{proof}
For an odd source choose the stage in Lemma~\ref{co:operators} with
cutoff $R\ge2$; only these fixed divisibilities and finitely many
bounded-weight coefficient tests are needed below.
Two fresh traces in matching parity, at least one split, add weight
at least $3/2$.  The allowed valuation in \eqref{co:BU-bound} is an
integer, so it is at least $2$ before the unary division and the
result vanishes after it.  In opposite parity any two fresh traces
add weight at least $1$; subtracting $e_U=1/2$ and using integral
valuation again gives vanishing.  The estimates persist under
further tests.  Hence respectively
$t(g)t(u)=0$ for split $u$, or $t(g)t(u)=0$ for every $u$.
Applying \eqref{co:trace-identities} to $t(gu)t(u)$ gives
$t(gu^2)=t(g)$.  The quotient by the subgroup generated by squares
is elementary abelian, and the Galois systems are unramified off
their level.  This proves the stated factorization.  At coefficient
$q$ the fresh split weight $1$ pays exactly the division in $H_U$;
the fresh simple weight $1/2$ pays exactly $e_U$ in opposite parity.
No unary term survives.  This proves the symbol-reading assertion.

For contraction, apply \eqref{co:inertia-identity} at coefficient
$1$, so that its left side is a coefficient-$p$ test.  Replace its
trace arguments by fresh primes.  The first-shell calculation in
Lemma~\ref{co:operators} gives, in each of the split--split,
split--simple, and simple--split cases,
\[
 c_p(D')A(D'),\qquad D'=(U/p)q',\quad
 \operatorname{Fr}_{q'}\sim g\phi_p.
\]
Here the omitted source is exactly $B_V$ or $H_V$, with
$V=U\setminus\{p\}$, constructed from the same original modular form
and the same $\eta,\psi,\delta$;
no new minimum or normalization is chosen.  To check the integer
powers explicitly, put $n_U=\lceil\delta+w(U)\rceil$ and let
$m_U$ be $1$ in matching parity and $0$ otherwise.  After fresh
tests kill the unary terms, the first shell relative to the
omitted-source test has multiplier
\[
 2^{1+n_V-n_U-m_U+m_V}c_p(D').
\]
For split $p$ the exponent is zero.  For simple $p$ and matching
$U$ it is also zero.  For simple $p$ and opposite $U$ it is one;
this is precisely the simple--simple case.  The character change
$\psi_U/\psi_V=(p/\cdot)$ is already the Legendre factor in this
shell calculation.  The phase is a unit and disappears on reduction.
Let $q_0',q_1'$ realize $g\phi_p,g\tau_p\phi_p$, respectively,
and put $D_i'=Vq_i'$.  Applying inertia changes the Legendre term
in \eqref{co:cp} by a sign, so $c_p(D_1')=c_p(D_0')+1$.
Although these Frobenius elements may first be specified in the full
current elementary extension, their restrictions to the omitted-source
extension agree: inertia
at either place above $p$ is killed in that extension.  Its original
characters are unramified at $p$, and its only new character factors
come from $V$.  Thus it is unramified at the whole rational prime
$p$, giving $A(D_1')=A(D_0')$.  If $\Delta_{\tau_p}A(Uq)$ denotes
the difference of the original tests with Frobenius labels $g\tau_p$
and $g$, the reduced local identity is now explicitly
\[
 \Delta_{\tau_p}A(Uq)
   =c_p(D_1')A(D_1')+c_p(D_0')A(D_0')=A(D_0'),
\]
which is \eqref{co:contraction}.  When both vertices are simple the product
Frobenius is split and the first-shell valuation has a positive
excess; its reduction is zero.  The entire-prime assertion follows
from the actual support of the omitted-$p$ coefficient form, rather
than from any convention about a chosen orientation.  Isolation
is the same symbol-reading test with $t(1)=0$.

For the final assertion first use two copies of a single simple
Frobenius trace on matching $H_U$.  Put
$b=B_U[1]$, $\epsilon=\chi_{-4}(q)$, and
$\kappa=(\psi_U/\eta)(q)$.  Substitution of
\eqref{co:recurrence} and of the unary character gives the two terms
in the coefficient-$1$ test of $T_q^2$, modulo $\mathfrak m$:
\begin{align*}
 H_U[q^2]&=\frac{b\eta(q)}2
        (a_q-\kappa\epsilon-q\epsilon^i),\\
 q(\langle q\rangle H_U)[1]&=\frac{b\eta(q)}2
        q(\kappa-\epsilon^{i+1}).
\end{align*}
Their sum is
\[
 \frac{b\eta(q)}2
 \left(a_q+\kappa(q-\epsilon)
                    -q\epsilon^i(1+\epsilon)\right).
\]
For $q\equiv3\pmod4$ the parenthesized half has residue $a_q/2$;
for $q\equiv1\pmod4$ it has residue $a_q/2+1$.
In both cases it is $\#E(\mathbb F_q)/2$ modulo $2$.
Since $b$ reduces to $A(U)$, the squared trace has the required
coefficient.  Fresh-prime testing replaces its two copies by
distinct primes.  To pass to any pair with the specified labels,
vary one prime at a time.  After that simple prime is omitted the
opposite-parity symbol depends only on rational elementary labels;
simple--simple changes vanish by the proved contraction rule.
Thus only the stated support and split-vertex labels matter.
\end{proof}
